\documentstyle[%


righttag%


,11pt%


,amssymb%


,russian%               remove in Eng.


,izvru%                 Set izven% in Eng.


]{amsart}





%





\newcommand{\diam}{\operatorname{diam}}


\newcommand{\tts}[1]{}





%\hoffset=-15mm%                  For Eng.


\setcounter{page}{55}


\god{2001}


\nomer{3~(466)} %                        Remove in Eng.


%\nomer{Vol.\,45, No.\,3}%               Set in Eng.


%\str{\pageref{begin}--\pageref{end}}%  Set in Eng.


\udk{515.124.4:514.774}





\author{\it Е.Н.\,Сосов}


% NB:   ^^ remove in Eng.


\title{О непрерывности и связности метрической $\delta$-проекции\\


в равномерно выпуклом геодезическом пространстве}


\thanks{Работа выполнена при финансовой поддержке Российского фонда


фундаментальных исследований, грант \No\,00-01-00308.}








\date{}


%\pagestyle{myheadings}%         Set in Eng.


\pagestyle{plain} %              Remove in Eng.


\begin{document}


%\thispagestyle{plain}%          Set in Eng.


\maketitle


\label{begin}





В статье некоторые результаты работ [1]--[3]


о непрерывности и связности метрической $\delta $-проекции в равномерно


выпуклом банаховом пространстве обобщаются на случай равномерно


выпуклого геодезического пространства. Одним из простых следствий такого


обобщения является справедливость аналогичных результатов в пространствах


Лобачевского (включая бесконечномерные). 





\section{Необходимые определения и теоремы}





Пусть $(X,\rho )$ --- выпуклое метрическое пространство [3]


(т.\,е. метрическое пространство,


в котором расстояние между каждыми двумя точками равно длине кратчайшей


кривой с концами в этих точках).


Выпуклое метрическое пространство назовем прямым выпуклым 


метрическим пространством (ср. [4], с.\,13), если через каждые его две 


различные точки можно провести единственную прямую


(т.\,е. геодезическую кривую, изометричную всей вещественной 


оси $R$ со стандартной метрикой ([4], с.\,52)). Полное 


выпуклое метрическое пространство называется 


геодезическим пространством [5].





В дальнейшем используем следующие обозначения:


$B(x,r)$ ($B[x,r]$, $S(x,r)$) --- открытый шар (замкнутый шар, сфера) с


центром в точке $x$, радиуса $r>0$, $xy = \rho (x,y)$, $xM = \rho (x,M)$;


$[x,y]$ ($(x,y)$) --- замкнутый (открытый) отрезок с концами $x, y \in X$.


Пусть $\omega_{\lambda}(x,y)$, $\lambda \in R$, --- точка на 


прямой прямого выпуклого метрического пространства, 


проходящей через точки $x$, $y$, удовлетворяющая условиям 


а)~$x\omega_{\lambda}(x,y) = |\lambda|xy$; б)~если $\lambda \in [0,1]$,


то $\omega_{\lambda}(x,y) \in [x,y]$; в)~если $\lambda > 1$, то $y \in


(x,\omega_{\lambda}(x,y))$; г)~если $\lambda < 0$, то $x \in


(\omega_{\lambda}(x,y),y)$.


 


Множество $M$ прямого выпуклого метрического пространства


называется выпуклым, если для каждых двух различных точек $x$, $y$


из этого множества отрезок $[x,y]$ 


принадлежит этому множеству.


\medskip





{\bf Определение.}


Прямое выпуклое метрическое пространство $X$ назовем равномерно 


выпуклым метрическим пространством, если выполняются следующие условия:


\begin{itemize} 


\item[A)] для каждого $\lambda \in R$ отображение $\omega_{\lambda}:


X\times X\to X$ равномерно непрерывно на каждом множестве вида $B\times B$,


где $B$ --- произвольный замкнутый шар пространства $X$;


\item[B)] каждый замкнутый шар $B[p,r]$ пространства $X$ выпуклый и обладает


свойством: если $\lim \limits_{n \to \infty}{p\omega_{1/2}(x_{n},y_{n})}=r$


для последовательностей $(x_{n}), (y_{n})$ из шара $B[p,r]$, то ([6])


$$


\lim \limits_{n \to \infty}{x_{n}y_{n}} = 0.


$$


\end{itemize}


\medskip





Простыми примерами равномерно выпуклых геодезических пространств являются


равномерно выпуклые банаховы пространства и пространства Лобачевского 


(включая бесконечномерные).





Напомним следующие обозначения и определения из [2], [3]. Пусть


$\delta \geq 0$, $x \in X$, $M \subset X$, $M \neq \emptyset $, $t \geq 0$;


$x_{M}^{\delta} = \{y \in M : xy \le xM + \delta \}$


($ x_{M} = x_{M}^{0}$) --- метрическая


$\delta $-проекция (метрическая проекция) точки $x$ на множество $M$;


$P : X\times 2^{X}\times R_{+} \rightarrow 2^{X}$,


$P(x,M,\delta ) = x_{M}^{\delta}$


--- оператор метрического $\delta $-проектирования (оператор метрического 


проектирования, если $\delta \equiv 0$);


$\lambda (M) = \sup \{\rho (A,B) : A\cup B = M$, $A,B \neq \emptyset \}$;


$\lambda _{M}(\delta ) = \sup \{\lambda (x_{M}^{\delta}) : x \in X \}$


--- мера несвязности метрической $\delta$-проекции;


$\beta (A,B) = \sup \{xB : x \in A \}$,


$\alpha (A,B) = \max \{\beta (A,B),\beta (B,A)\}$ ---


полууклонение и расстояние по Хаусдорфу между непустыми, ограниченными


множествами $A, B \subset X$;


$\omega (t) = \sup \{\alpha (x_{M}^{\delta},y_{M}^{\delta}):y\in X$,


$xy \leq t, \ y_{M}^{\delta} \neq \emptyset \}$,


$\wh {\omega}(t) =


\sup \{\alpha (x_{M}^{\delta},x_{N}^{\delta}) : N \subset X$, 


$\alpha (M,N) \leq t, \ x_{N}^{\delta} \neq \emptyset \}$,


$\ov {\omega}(t) = \sup \{\alpha (x_{M}^{\delta},x_{M}^{\varepsilon }):


\varepsilon \geq 0$,


$|\varepsilon - \delta | \leq t,\ x_{M}^{\varepsilon} \neq \emptyset \}$,


$\Omega (t) = \sup \{\alpha (x_{M}^{\delta},y_{N}^{\varepsilon }) :


\varepsilon \geq 0$, 


$y \in X$, $N \subset X$, $xy \leq t$,


$|\varepsilon - \delta | \leq t$, $\alpha (M,N) \leq t$,


$y_{N}^{\varepsilon } \neq \emptyset \}$ --- четыре модуля


непрерывности оператора $P$ в точке $(x,M,\delta)$


по $x$, по $M$, по $\delta $ и по


совокупности этих переменных, причем $x_{M}\neq \emptyset$ при $\delta=0$;


$u_{M}(\delta ) = \sup \{\omega (+0) : x \in X \}$, $\wh {u}_{M}(\delta ) = 


\sup \{\wh {\omega }(+0) : x \in X \}$,


$\ov {u}_{M}(\delta ) = \sup \{\ov {\omega}(+0) : x \in X \}$,


$U_{M}(\delta ) = \sup \{\Omega (+0) : x \in X \}$,


где


$\omega (+0) = \lim \limits_{t \rightarrow +0}{\omega (t)}$


--- верхние грани разрывов


оператора $P$ на множестве $\{(x,M,\delta ) : x \in X \}$.


Оператор $P_{M} : X \rightarrow 2^{X}$,


$P_{M}(x) = x_{M}$ метрического проектирования 


на множество $M$ называется непрерывным в слабом смысле,


если $x_{n} \rightarrow x$ 


влечет $\rho (P_{M}(x_{n}),P_{M}(x)) \rightarrow 0$ [1]. 


$E_{M} = \{ x \in X : x_{M} \neq \emptyset \}$,


$T_{M} = \{ x \in X : x_{M}$ одноточечно$\}$,


$D_{M}x = \lim \limits_{\delta \to +0}{\diam(x_{M}^{\delta})}$,


$T'_{M} = \{ x \in X : D_{M}x = 0 \}$.


Множество $M \subset X$ называется $P$-связным, если для каждого $x \in X$ 


множество $x_{M}$ непустое и связное. Множество $M \subset X$ называется


$B$-связным ($\overset{\circ}{B}$-связным),


если пересечение множества $M$ с каждым 


замкнутым (открытым) шаром является связным. 


Множество $M \subset X$ называется множеством существования (чебышевским


множеством), если $E_{M} = X$ ($T_{M} = X$).





Приведем несколько вспомогательных утверждений.





\begin{lem}[{\series{m}\shape{n}\selectfont [6]}]


Пусть $X$ --- равномерно выпуклое метрическое пространство и $x \in X$. Если


последовательности $(y_{n})$, $(z_{n})$, $(u_{n})$ удовлетворяют


условиям $xy_{n} = r$ $(r>0)$, $xz_{n} = r + h$ $(h>0)$, $u_{n} \in


[x,z_{n}]$, $xu_{n} = r$ $(n = 1,2,\dotsc)$,


$\lim \limits_{n \to \infty}{y_{n}z_{n}} =h$, 


то $\lim \limits_{n \to \infty}{y_{n}u_{n}} = 0$.


\end{lem}





\begin{prop}[{\series{m}\shape{n}\selectfont [6], теорема 2}]


Пусть $X$ --- равномерно выпуклое геодезическое пространство,


а $M$ --- непустое замкнутое множество в пространстве $X$. Тогда каждое из


множеств $T_{M}$, $T'_{M}$ является дополнением множества первой


категории $($в частности, всюду плотно$)$.


\end{prop}





\begin{prop}[{\series{m}\shape{n}\selectfont [3], теорема 1}]


Пусть $M$ --- подмножество выпуклого метрического пространства $X$. 


Тогда $\lambda_{M}(\delta + 0) \leq u_{M}(\delta )$


при $0 < \delta < \lambda _{M}(\delta + 0)$,


а если $M$ --- множество существования, то справедливы неравенства 


$\lambda _{M}(\delta ) \leq u_{M}(\delta )$


при $0 \leq \delta < \lambda _{M}(\delta );$


$\max [\lambda _{M}(+0),\lambda _{M}(0)] \leq u_{M}(0)$.


\end{prop}





Сформулируем теперь полученные результаты. 





\begin{thm}


Пусть множество $M$ принадлежит прямому геодезическому пространству $X$, 


в котором каждый замкнутый шар выпуклый и, кроме того, выполняется одно из


следующих условий\/$:$


\begin{itemize}


\item[a)] множество $M$\; $P$-связно и оператор


$P_{M}$ непрерывен в слабом смысле$;$


\item[b)] $\lim \limits_{\delta \to +0}


{\lambda (x_{M}^{ \delta})} = 0$ для каждого $x \in X$.


\end{itemize}


Тогда $M$ --- $\overset{\circ}{B}$-связное множество. 


\end{thm}





\begin{lem}


Пусть $X$ --- равномерно выпуклое метрическое пространство,


$x \in X$, $0 < h <H$.


Тогда {\em a)}~$\diam(B[z,zy + \delta]\setminus B(x,xy)) \to 0;$ 


{\em b)}~$\diam(B[z,zy]\setminus B(x,xy - \delta)) \to 0$


при $\delta \to +0$ равномерно по 


всем $y, z \in X$, удовлетворяющим условиям


$z \in (x,y)$, $xy = H$, $xz = h$.


\end{lem}





\begin{thm}


Пусть $M$ --- $P$-связное множество в равномерно выпуклом геодезическом 


пространстве $X$. Тогда $M$ --- $B$-связное множество.


\end{thm}





\begin{lem}


Пусть $M$ --- замкнутое множество в равномерно выпуклом геодезическом 


пространстве $X$. Тогда следующие утверждения эквивалентны$:$





{\em a)} множество $M$\; $\overset{\circ}{B}$-связно$;\quad$


{\em b)} $x_{M}^{ \delta}$ связно для каждого $x \in X$


и каждого $\delta > 0;$ 





{\em c)} $\lambda _{M}(\delta ) = 0$


для каждого $\delta > 0;\quad$


{\em d)} $\lim \limits_{\delta \to +0}{\lambda (x_{M}^{ \delta})} = 0$


для каждого $x \in X$.


\end{lem}





\begin{lem}


Пусть $X$ --- прямое выпуклое метрическое пространство.


Тогда для каждой точки 


$(x,M,\delta) \subset X\times 2^{X}\times R_{+}$,


$\delta > 0$, $M \neq \emptyset$,


выполняется неравенство


$\ov {\omega}(t) < \wh {\omega}(t) + t$ при $0 < t < \delta $.


\end{lem}





\begin{note}


Из леммы 4 и неравенства (2) из [3]:


$\Omega (t) \leq \ov {\omega}(t) + t$


для $(x,M,\delta) \subset X\times 2^{X}\times R_{+}$,


$\delta > 0$, $M \neq \emptyset$,


в прямом выпуклом метрическом пространстве $X$


следуют соотношения, аналогичные соотношениям (4) в [3],


\begin{equation}


u_{M}(\delta ) \leq \wh {u}_{M}(\delta ) =


\ov {u}_{M}(\delta )= U_{M}(\delta ) 


\ \; \text{для каждого} \ \; \delta > 0.


\end{equation}


\end{note}





\begin{thm}


Пусть $X$ --- равномерно выпуклое метрическое пространство,


$M \subset X$, $\delta > 0$.


Тогда $U_{M}(\delta ) \leq \lambda (x_{M}^{ \delta - 0})$.


\end{thm}





\begin{thm}


Пусть $X$ --- равномерно выпуклое метрическое пространство, $M\subset X$.


Тогда нижеследующие условия {\em a)--e)} эквивалентны.


Если же $M$ --- замкнутое 


множество равномерно выпуклого геодезического пространства $X$,


то условия {\em a)--h)} эквивалентны.





{\em а)} $U_{M}(\delta ) = 0$ для каждого $\delta > 0;$


{\em b)} $\lim \limits_{\overline{\delta \to +0}}{U_{M}(\delta )} = 0;$


 


{\em c)} $u_{M}(\delta ) = 0$ для каждого $\delta > 0;$ 


{\em d)} $\lim \limits_{\overline{\delta \to +0}}{u_{M}(\delta )} = 0;$





{\em e)} $\lambda _{M}(\delta) = 0$ для каждого $\delta > 0;$ 


{\em f)} $\lim \limits_{\delta \to +0}{\lambda (x_{M}^{ \delta})} = 0$


для каждого $x \in X;$





{\em g)} $M$ --- $\overset{\circ}{B}$-связное множество$;$ 


{\em h)} $x_{M}^{ \delta}$ --- связное множество для каждых


$\delta > 0$ и $x \in X$.


\end{thm}





\begin{cor}[{\series{m}\shape{n}\selectfont теорем 2, 3}]


Пусть $M$ --- чебышевское множество в равномерно выпуклом геодезическом


пространстве $X$. Тогда оператор $P$ метрического $\delta$-проектирования на


множество $M$ непрерывен по совокупности переменных,


т.\,е. $U_{M}(\delta ) = 0$


для каждого $\delta > 0$.


\end{cor}





\begin{cor}[{\series{m}\shape{n}\selectfont предложения 2 и теоремы 3}]


Пусть $M$ --- множество существования из равномерно выпуклого метрического


пространства $X$, в котором оператор $P_{M}$ метрического проектирования на


множество $M$ непрерывен по Хаусдорфу (т.\,е. $u_{M}(0) = 0$). Тогда 


$U_{M}(\delta ) = 0$ для каждого $\delta > 0$.


\end{cor}





\begin{note}


Теорема 1 является обобщением теоремы 1 Л.П.\,Власова [1] и леммы 2 


А.В.\,Маринова [3]. Теорема 2 является обобщением


теоремы 4.2 Л.П.\,Власова [2].


Теоремы 3, 4 являются обобщениями теорем 2, 4 А.В.\,Маринова [3].


Утверждение~a) леммы 2 является обобщением леммы 1.1 С.Б.\,Стечкина 


[1] и доказано в [6], утверждение~b) леммы 2


является обобщением утверждения в лемме 4


А.В.\,Маринова [3] и доказано в этой статье. Леммы 3, 4 являются обобщениями 


леммы 3 и неравенства (3) А.В.\,Маринова [3].


\end{note}





\section{Доказательства полученных результатов}





{\bf Доказательство теоремы 1}.


Доказательство при условии a) почти дословно повторяет


доказательство теоремы 1 Л.П.\,Власова [1].


Отличие --- в содержании понятия деления отрезка в данном отношении, 


своевременном упоминании условия выпуклости каждого 


замкнутого шара и обозначениях.


Пусть, напротив, найдется такой открытый шар


$B_{0} = B(x,R_{0})$, что множество 


$M\cap B_{0}$ несвязно. Значит, $M\cap B_{0} = A\cup C$,


где $A$, $C$ --- непустые, замкнутые в $B_{0}$ 


и непересекающиеся множества. Выберем $R_{1}$ так, чтобы 


$\max [xA,xC] < R_{1} < R_{0}$.


Тогда замкнутый шар $B_{1} = B[x,R_{1}]$ обладает тем 


свойством, что $B_{1}\cap M$ распадается на непустые, замкнутые и 


непересекающиеся множества $A\cap B_{1}$, $C\cap B_{1}$. 


Покажем, что найдется замкнутый шар $B = B[y,r]\subset B_{0}$ такой, что 


$\rho (B\cap A,B\cap C) > 0$.


Допустим, что это не так. Пусть $R_{i} > 0$ при $i \geq 2$


такие, что $\sum \limits_{i=2}^{\infty }{R_{i}} < R_{0} - R_{1}$. Построим


последовательность замкнутых шаров


$B_{1}, \, B_{2}, \dotsc, B_{n}, \dotsc$ индуктивно.


Шар $B_{1}$ уже построен. Пусть


$B_{n}\cap A \neq \emptyset$, $B_{n}\cap C \neq \emptyset$ 


и радиус $B_{n}$ равен $R_{n}$. Тогда по предположению 


$\rho (B_{n}\cap A,B_{n}\cap C) = 0$.


Значит, найдутся точки $a_{n} \in B_{n}\cap A$,


$c_{n} \in B_{n}\cap C$ такие, что $a_{n}c_{n} < 2R_{n+1}$.


В качестве $B_{n+1}$


возьмем шар $B[q_{n+1},R_{n+1}]$,


где $q_{n+1} = \omega _{1/2}(a_{n},c_{n})$. Из условия 


выпуклости каждого замкнутого шара и неравенства треугольника следует,


что $q_{n}q_{n+p} \leq q_{n}q_{n+1} + \dotsb + q_{n+p-1}q_{n+p} \leq 


\sum \limits_{k=n}^{n+p-1}{R_{k}} \rightarrow 0$


при $p\rightarrow \infty$, $n\rightarrow \infty$.


Значит, $(q_{n})$ --- фундаментальная 


последовательность и найдется такая точка 


$q \in X$, что $q_{n}\rightarrow q$ ($n\rightarrow \infty$).


Кроме того, $x = q_{1}$ и


$xq = q_{1}q \leq q_{1}q_{n} +


q_{n}q \leq \sum \limits_{k=1}^{n-1}{q_{k}q_{k+1}} + q_{n}q


\leq \sum \limits_{k=1}^{\infty }{R_{k}} + q_{n}q$.


Но $\sum \limits_{k=1}^{\infty }{R_{k}} < R_{0}$, $q_{n}q\rightarrow 0$


($n\rightarrow \infty$).


Следовательно, $q \in B_{0}$. Кроме того, $R_{n}\rightarrow 0$


($n\rightarrow \infty$).


Значит, $a_{n}\rightarrow q$, $c_{n}\rightarrow q$ ($n\rightarrow 


\infty$) и $q \in A\cap C\cap B_{0}$.


Противоречие. Таким образом, вышеуказанный шар


существует. Тогда $B\cap M = A'\cup C'$, где $A' = A\cap B$,


$C' = C\cap B$ --- непустые


замкнутые множества с $\rho (A',C') > 0$.





1)~Значит, $y_{M} \subset A'\cup C'$.


Пусть для определенности $y_{M}\subset A'$. Тогда


$y_{M} \not \subset C'$, поскольку $\rho (A',C') > 0$


и $y_{M}$ связно. Пусть $c \in C'$.


Тогда $P_{M}(\omega_{\lambda }(y,c)) \subset B$


при $0 \leq \lambda \leq 1$, т.\,к.


$\omega _{\lambda }(y,c)M \leq \omega _{\lambda }(y,c)c


\leq R - y\omega _{\lambda }(y,c) =


\rho (\omega _{\lambda }(y,c),X\setminus B)$.


Значит, $P_{M}(\omega_{\lambda }(y,c))


\subset A'\cup C'$. Выберем $\lambda _{0} =


\sup \{\lambda : 0 \leq \lambda \leq 1$,


$P_{M}(\omega_{\lambda }(y,c)) \subset A' \}$.


Пусть последовательность $(\lambda _{n})$


($n=1, 2,\dotsc$) такая, что


$\lambda _{n}\rightarrow \lambda _{0}$, $\lambda _{n} > \lambda _{0}$.


Тогда $P_{M}(\omega_{\lambda_{n} }(y,c)) \subset C'$,


$\omega_{\lambda_{n} }(y,c)\rightarrow 


\omega_{\lambda_{0} }(y,c)$


и $\rho (P_{M}(\omega_{\lambda_{n} }(y,c)),


P_{M}(\omega_{\lambda_{0} }(y,c)))\rightarrow 0$ ($n\rightarrow \infty$).


Следовательно, $\rho (A',P_{M}(\omega_{\lambda_{0} }(y,c))) = 0$.


Но возможны только два случая:


$P_{M}(\omega_{\lambda_{0} }(y,c)) \subset A'$


или $P_{M}(\omega_{\lambda_{0} }(y,c))


\subset C'$. Значит, $\rho (A',C') = 0$. Противоречие. 


Теорема 1 при условии а) доказана.





2) (Cр. доказательство леммы 2 в [3].) Если множество $M$ не является


$\overset{\circ}{B}$-связным, то, как было доказано выше, найдутся


$x \in X$ и $\delta > 0$ такие, что $x_{M}^{ \delta} = A'\cup C'$,


$\rho (A',C') > 0$, $\max [xA',xC'] < xM + \delta $.


Нетрудно заметить, что в случае $xA' = xC'$ верно неравенство


$\rho (A',C') \leq \lambda (x_{M}^{ \delta})$ при 


$\delta > 0$. Но это приводит к противоречию с условием~b) теоремы.


Для завершения


доказательства осталось привести случай $xA' < xC'$ к предыдущему.


Пусть $0 < 2t < xM + \delta - xC'$, $y \in x_{C'}^{t}$.


Найдется такая точка $v \in [x,y]$,


что $vA' = vC'$, т.\,к. непрерывная функция


$zA' - zC'$ меняет знак на отрезке $[x,y]$. Тогда


$xv = xy - vy \leq xC' + t - vC'$


и $B[v,vC' + t] \subset B[x,xM + \delta ]$. Значит,


$v_{M}^{t} \subset A'\cup C'$


и $\rho (A',C') \leq \lambda (v_{M}^{t})$. А это противоречит


условию~b).$\quad\square$


\medskip





{\bf Доказательство п.\,b) леммы 2}. Пусть для каждого $0<\delta<H-h$ в


множестве $B[z,zy]\setminus B(x,xy - \delta)$


произвольным образом выбирается элемент $u$.


Очевидно, достаточно доказать, что $uy \to 0$ при $\delta \to +0$.


Пусть $v$ --- точка пересечения прямой,


проходящей через точки $x$, $u$, со сферой


$S(x,xy)$, $t = [x,u]\bigcap S(x,xz)$. Из неравенств


$H - \delta \leq xu \leq xv = H$ следует,


что $xu \rightarrow H$ и $uv =H - xu \rightarrow 0$


при $\delta \to +0$. Тогда из неравенств


$zy \leq zv \leq zu + uv \leq zy + uv$ следует, что


$zv \rightarrow H - h$ при $\delta \to +0$.


Но $xt = xz = h$. Поэтому по лемме 1 имеем $tz \to 0$


при $\delta \to +0$. Из условия A) определения


следует теперь, что $yv \to 0$ при


$\delta \to +0$. Таким образом, $uy \leq uv +vy \to 0$


при $\delta \to +0$.$\quad\square$


\medskip





{\bf Доказательство теоремы 2} получается из доказательства теоремы 4.2 


Л.П.\,Власова [2] заменой ссылок на лемму 1.1 и теорему


1.1 С.Б.\,Стечкина [2] ссылками


на их обобщения: лемму 2 и предложение 1 соответственно, а также заменой 


обозначений.$\quad\square$


\medskip





{\bf Доказательство леммы 3}. Импликации


b)\,$\Rightarrow\,$\,c)\,$\Rightarrow$\,d)


очевидны.


Импликация d)\,$\Rightarrow\,$a) доказана в п.~2) доказательства


теоремы~1. Импликация


a)\,$\Rightarrow\,$b) отличается от доказательства в лемме 3


А.В.\,Ма\-ри\-но\-ва только тем,


что ссылка на доказательство теоремы 4.2 Л.П.\,Власова должна быть заменена


ссылкой на доказательство теоремы~2.$\quad\square$





{\bf Доказательство леммы 4}.


Поставим в соответствие каждой точке $z \in x_{M}^{t}$ 


точку $v(z) = S(x,xM +t)\cap \{\omega _{l}(x,z) : l \geq 0 \}$. Положим 


$N = \{v(z) : z \in x_{M}^{t} \} \cup (M\setminus x_{M}^{t})$. Тогда 


$\alpha (M,N) = \beta (x_{M}^{t + \delta},x_{N}^{\delta }) = t$.


Следовательно,


$\beta (x_{M}^{t + \delta},x_{M}^{\delta }) \leq


\beta (x_{M}^{t + \delta},x_{N}^{\delta }) +


\beta (x_{N}^{\delta},x_{M}^{\delta }) \leq t + \wh {\omega}(t)$.





Положим теперь


$W = \{\omega _{l}(x,z) : z \in x_{M}^{\delta }


\setminus x_{M}^{\delta - t}$, $l\rho (x,z) = xz + t \}$,


$K = [M\setminus (x_{M}^{\delta } \setminus x_{M}^{\delta - t})]\cup W$.


Тогда 


$W\cap x_{M}^{\delta } = \emptyset $,


$x_{K}^{\delta } = x_{M}^{\delta - t} $, 


$\alpha (M,K) \leq t$. Отсюда получим 


$\beta (x_{M}^{\delta },x_{M}^{\delta - t})


\leq \beta (x_{M}^{\delta },x_{K}^{\delta }) +


\beta (x_{K}^{\delta },x_{M}^{\delta - t}) =


\beta (x_{M}^{\delta },x_{K}^{\delta}) \leq \wh {\omega}(t)$.$\quad\square$


\medskip





{\bf Доказательство теоремы 3} получается из доказательства теоремы 2 


А.В.\,Ма\-ри\-но\-ва [3] заменой ссылки на лемму 4 в [3] ссылкой на 


более общую лемму 2 данной статьи и заменой обозначений.$\quad\square$


\medskip





{\bf Доказательство теоремы 4} (ср. [3]).


Из предложения 2 и теоремы 3 следует


справедливость импликаций d)\,$\Rightarrow$\,e)\,$\Rightarrow$\,a).


Но из первых четырех


условий первое самое сильное, а последнее самое слабое. Значит, условия 


a)-- e) эквивалентны. Эквивалентность условий e)--h) при дополнительном 


предположении полноты пространства $X$ и замкнутости множества $M$ следует


теперь из леммы 3.$\quad\square$





\begin{thebibliography}{9}





\bibitem{1}


Власов Л.П. {\em Чебышевские множества и некоторые их обобщения}


// Матем. заметки.


-- 1968. -- Т.\,3. -- \No\,1. -- С.\,59--69.


\bibitem{2}


Власов Л.П. {\em Аппроксимативные свойства множеств в линейных


нормированных пространствах} // УМН. -- 1973. -- Т.\,28. --


Вып.\,6. -- С.\,3--66.


\bibitem{3}


Маринов А.В. {\em Непрерывность и связность метрической


$\delta$-проекции} // Аппроксимация в конкрет. и 


абстракт. банахов. пространствах. -- Свердловск,


1987. -- С.\,82--95.


\bibitem{4}


Буземан Г. {\em Геометрия геодезических}. -- М.: Физматгиз, 1962. --


503~с.


\bibitem{5}


Ефремович В.А. {\em Неэквиморфность пространств Эвклида и Лобачевского} //


УМН. -- 1949. -- Т.\,4. -- Вып.\,2. -- С.\,178--179.


\bibitem{6}


Сосов Е.Н. {\em Об аппроксимативных свойствах множеств в специальных


метрических пространствах} // Изв. вузов. Математика. -- 1999.


-- \No\,6. -- С.\,81--84.





\end{thebibliography}





\vskip 2cm





\post{\it Казанский государственный}{\it Поступила}


\post{\it педагогический университет}{01.10.1999}





\label{end}


\end{document}








